%
% $Id: slides.tex 4228 2006-06-21 21:55:12Z jjamor $
%
%
% Compilar a .pdf con LaTeX (pdflatex)
% Es necesario instalar Beamer (paquete latex-beamer en Debian)
%

%
% Gr�ficos:
% Los gr�ficos pueden suministrarse en PNG, JPG, TIF, PDF, MPS
% Los EPS deben convertirse a PDF (usar epstopdf)
%

\documentclass{beamer}
\usetheme{Warsaw}
\usebackgroundtemplate{\includegraphics[width=\paperwidth]{format/libresoft-bg.png}}
\usepackage[spanish]{babel}
\usepackage[latin1]{inputenc}
\usepackage{graphics}
\usepackage{amssymb} % Simbolos matematicos
\usepackage{url}

%\definecolor{libresoftgreen}{RGB}{162,190,43}
%\definecolor{libresoftblue}{RGB}{0,98,143}

%\setbeamercolor{titlelike}{bg=libresoftgreen}

%% Metadatos del PDF.
\hypersetup{
  pdftitle={Libre Software Communities},
  pdfauthor={Daniel Izquierdo Cort�zar},
  pdfcreator={GSyC/Libresoft},
  pdfproducer=PDFLaTeX,
  pdfsubject={nn},
}
%%


\AtBeginSection[]
{
  \begin{frame}<presentation>
    \frametitle{Index}
    \tableofcontents[current]
  \end{frame}
}


\begin{document}

\title{Libre Software Communities: An Approach}
\subtitle{Master on Free Software}
\institute{dizquierdo@gsyc.escet.urjc.es\\
GSyC/Libresoft}
\author{Daniel Izquierdo Cort�zar}
\date{\today}

\frame{
\maketitle
\begin{center}
\includegraphics[width=6cm]{format/gsyc-urjc}
\end{center}
}


% Si el titulo o el autor se quieren acortar para los pies de p�gina
% se pueden redefinir aqu�:
%\title{Titulo corto}
%\author{Autores abreviado}


%% LICENCIA DE REDISTRIBUCION DE LAS TRANSPAS
\frame{
~
\vspace{4cm}

\begin{flushright}
{\tiny
(cc) 2007 Daniel Izquierdo Cort�zar. \\
Some rights reserved. This document is distributed under the Creative \\
            Commons Attribution-ShareAlike 2.5 licence, available in \\
            http://creativecommons.org/licenses/by-sa/2.5/

%  Este documento (o uno muy similar) está disponible en \\
%  \url{http://gsyc.escet.urjc.es/~jjamor/}
}
\end{flushright}
}
%%



\section{Introduction}

\begin{frame}
\frametitle{Take a look around you}
\begin{itemize}
\item Is somebody doing the same work?
\item Try to look for information at:\\
Forges: SF - \url{http://sourceforge.net}, Berlios - \url{http://berlios.de} \\
Project's indexes: Freshmeat \- \url{http://freshmeat.net}, FSF directory \- \url{http://directory.fsf.org} 
\item Is better to unify efforts?
\end{itemize}
\end{frame}


\begin{frame}
\frametitle{To unify efforts or not}
\begin{itemize}
\item GNOME - KDE
\item Beryl - Compiz
\item Several distributions (Debian, Suse, Fedora, Mandriva, Ubuntu, Linex, GuadaLinex, ...)
\end{itemize}
\end{frame}


\begin{frame}
\frametitle{Take a look around you}
\begin{itemize}
\item Have you contacted with them?
\item Different goals?, different points of view? just want to start a project? others ...
\item Let's go!
\end{itemize}
\end{frame}


\begin{frame}
\frametitle{First steps}
\begin{itemize}
\item Mission limits\\
What your product will do... (List of features)\\
What your product will not do
\item Name \\
\begin{tiny}
\begin{quotation}
What's your name? \\
{\em Apache!} \\
Nice to meet you :)
\end{quotation}
\end{tiny}
\item Have a clear mission statement
\item Alpha, Beta
\end{itemize}
\end{frame}

\section{Legal Aspects}

\begin{frame}
\frametitle{License, what is the best?}
\begin{itemize}
\item Source code (remember): \\
GPL family (+ restrictive)\\ 
BSD family (- restrictive)

\item Documentation: \\
Creative Commons? 

\item You could create your own license, however...\\


\end{itemize}
\end{frame}



\begin{frame}
\frametitle{License}
\begin{itemize}

\item It is usual to start from a minor project\\

\item Perhaps there are several authors: \\

A company. \\

%You could have legal problems if the source code was made by them and they disagree
%in the process of liberation.\\

%Sometimes you sign a contract like: ``All the things you create, even in your free time
%will have our copyright''.\\

%Some free software communities need a contract to confirm that the company wants to 
%free its software. It occurs when the community or foundation will get the copyright, thus
%they can get in trouble.\\

You are the owner.\\

Several authors worked on it.\\
% You have to ask them about changing the license, if not
%you could have legal problems.\\

\end{itemize}
\end{frame}


\begin{frame}
\frametitle{Pieces of advice}
\begin{itemize}
\item Your web site should contain a clear reference to the license you use.
\item Source code:\\
License could be included in each file.\\
Every file could contain a reference to the license (COPYING or LICENSE)
\end{itemize}
\end{frame}


\section{Image}

\subsection{Web site}

\begin{frame}
\frametitle{Visibility}
\begin{itemize}
\item Basic to:\\
Attract new users, developers.\\
More users = more bug reports = more activity = ...\\
\end{itemize}
\end{frame}


\begin{frame}
\frametitle{Structure}
\begin{itemize}
\item Basic information (Remember! first steps!)
\item Hot news
\item Activity and maintenance in mailing lists, wiki, trac ...
\item Related projects
\item Projects which use our product
\end{itemize}
\end{frame}



\begin{frame}
\frametitle{Documentation}
\begin{itemize}
\item User documentation
\item Developer documentation: specific mailing lists, SCM URL, BTS, others
\item Developer documentation: architecture, style, rules, others
\end{itemize}
\end{frame}


\begin{frame}
\frametitle{Examples}
\begin{itemize}
\item Plone
\item GNOME projects
\item SourceForge
\end{itemize}
\end{frame}


\subsection{Source code}

\begin{frame}
\frametitle{Users}
\begin{itemize}
\item Expert, medium and beginner\\
\item Easy, clear and fast installation\\
\item Good documentation\\
\end{itemize}
\end{frame}


\begin{frame}
\frametitle{Installation process}
\begin{itemize}
\item Difficulties mean:\\
Lost time answering doubts\\
Lost time fixing problems\\
Lost time updating the web site\\
\end{itemize}
\end{frame}


\begin{frame}
\frametitle{Installation process}
\begin{itemize}
\item If it fails, there are other products on the Internet!
\end{itemize}
\end{frame}

\begin{frame}
\frametitle{Cleaning code}
\begin{itemize}
\item Your code is your image (Yep, They are the developers of the product X)
\item Commits mailing list (everybody will know about your changes)
\end{itemize}
\end{frame}


\begin{frame}
\frametitle{Files}
\begin{itemize}
\item Source code must have several ``burocratic'' files: \\
License (LICENSE or COPYING)\\
General Info (README)\\
Changes made (Changelog) \\
Installation process info (INSTALL) \\
Documentation (Users, developers)\\
\end{itemize}
\end{frame}


\begin{frame}
\frametitle{Packages or binaries}
\begin{itemize}
\item They should be marked with the correct version
\item They must follow the same standard
\end{itemize}
\end{frame}


\begin{frame}
\frametitle{Other recomendations}
\begin{itemize}
\item Clean your code (use English)
\item Documentation (local language + English)
\item Try to prepare your software to be used by others\\
Translators: */po/*\\
Developers: Clear architecture and comments\\
...

\end{itemize}
\end{frame}



\begin{frame}
\frametitle{Demos}
\begin{itemize}
\item At least images with results (graphic environment)
\item or data in a database, XML, CSV, ...
\end{itemize}
\end{frame}

\section{Tools}
\begin{frame}
\frametitle{Web site}
\begin{itemize}
\item Collaborative development - Wiki
\item Fast installation - Zope/Plone
\item CMS in general
\end{itemize}
\end{frame}


\begin{frame}
\frametitle{}
\begin{itemize}
\item Each project has its goals:\\
\item SCM - SVN, CVS, Git ...
\item BTS - Bugzilla, Gnat, Trac...
\item Trac
\end{itemize}
\end{frame}


\section{Announcing}
\begin{frame}
\frametitle{Ok, but where?}
\begin{itemize}
\item SourceForge, Freshmeat, FSF directory
\item Your own blog
\item Basic scheme:\\
\begin{quotation}
What is it?\\
Where to get it?\\
Features\\
Requirements\\
Current version
\end{quotation}
\end{itemize}
\end{frame}



\section{Conclusions}

\begin{frame}
\frametitle{Conclusions}


\end{frame}



\end{document}
